\documentclass[Screen16to9,17pt]{foils}
\usepackage{zencurity-slides}
\externaldocument{software-security-exercises}
\selectlanguage{english}

\begin{document}

\mytitlepage
{11. Network Attacks}
{EK  Software Security}


\slide{Plan for today}

\begin{list1}
\item Subjects
\begin{list2}
\item Auditing Application Protocols
\item Example protocols and vulnerabilities
\item Abstract Syntax Notation (ASN.1) problems
\item Domain Name System (DNS) problems
\item Firewalls and related issues
\item Network Security Monitoring
\end{list2}
\item Exercises
\begin{list2}
\item Nmap and SYN flooding exercises
\end{list2}
\end{list1}

\slide{Reading Summary}

\begin{list1}
%\item \emph{Hacking, 2nd Edition: The Art of Exploitation}, Jon Erickson chapter 4, browse
\item Browse: \emph{TCP Synfloods - an old yet current problem, and improving pf's response to it}
Henning Brauer, BSDCan 2017 \link{http://quigon.bsws.de/papers/2017/bsdcan/}
\end{list1}



\slide{Goals: Introduction to Auditing Application Protocols}

\hlkimage{4cm}{wireshark-sni-twitter.png}

\begin{list1}
\item Often you dont need to audit the whole protocol in detail
\item Sometimes people can't tell which protocols, ports and services they use ...
\item And you need to configure a firewall/network filter
\item Picture: Wireshark with TLS SNI, recent Exim CVE-2019-15846 was SNI parsing
\item Plus SYN flooding, and Denial of Service
\end{list1}



\slide{What is pentesting and a pentest }

\begin{quote}
  A penetration test, colloquially known as a pen test, pentest or ethical hacking, is an authorized simulated cyberattack on a computer system, performed to evaluate the security of the system;[1][2] this is not to be confused with a vulnerability assessment.[3] The test is performed to identify weaknesses (also referred to as vulnerabilities), including the potential for unauthorized parties to gain access to the system's features and data,[4][5] as well as strengths,[6] enabling a full risk assessment to be completed.
\end{quote}
Source: quote from \link{https://en.wikipedia.org/wiki/Penetration_test}

\begin{list1}
\item Penetration testing is a simulation, with good intentions
\item People around the world constantly \emph{test your defenses}
\item Often better to test at planned times
\end{list1}

\slide{Time schedule}

\hlkimage{4cm}{kame-noanime-small.png}

Plan for the day
\begin{list1}
\item 12:45 - 13:45 -- Session 0: Teaching  Network Attacks intro
\item 13:45 - 14:00 -- 15min break
\item 14:00 - 14:45 -- Session 1: Exercise: *sshole
\item 14:45 - 15:00 -- 15min break
\item 15:00 - 16:45 -- Sessions 2+3: Teaching and exercises in Networks and DDoS
\end{list1}

\slide{Session 0: Teaching and exercise: Network Attacks intro}



\slide{Internet Today}

\begin{quote}
\textbf{\large Computer network}

In computer science, computer engineering, and telecommunications, \textbf{a network} is a \textbf{group} of communicating computers and peripherals called hosts, which communicate data to other hosts via communication protocols, as facilitated by networking hardware.

Within a computer network, \textbf{hosts} are identified by \textbf{network addresses}, which allow networking hardware to locate and identify hosts. Hosts may also have \textbf{hostnames}, memorable labels for the host nodes, which can be mapped to a network address using a hosts file or a name server such as \textbf{Domain Name Service}. The physical medium that supports information exchange includes \textbf{wired media like copper cables, optical fibers, and wireless radio-frequency media}. The arrangement of hosts and hardware within a network architecture is called the \textbf{network topology}.[1][2]
\end{quote}
Source: \url{https://en.wikipedia.org/wiki/Computer_network}



\slide{Internet Today}

\hlkimage{13cm}{images/server-client.pdf}

\begin{list1}
\item Clients and servers, roots in the academic world
\item Protocols are old, very old! RFC 791: STD 5: Internet Protocol is from  September 1981!
\item Very few protocols were encrypted, today a lot has switched to HTTPS and TLS
\end{list1}

\slide{Internet is Open Standards!}

{\hlkbig \color{titlecolor}
We reject kings, presidents, and voting.\\
We believe in rough consensus and running code.\\
-- The IETF credo Dave Clark, 1992.}

\begin{list1}
\item Request for comments (RFC) -- a series of documents spanning decades
\item RFC, BCP, FYI, informational -- first ones from 1969!
\item Are not updated but status is changed to Obsoleted when new versions are published
\item Standards track:\\
Proposed Standard $\rightarrow$ Draft Standard $\rightarrow$ Standard
\end{list1}


\slide{Network data at different layers}

\hlkimage{16cm}{sample-ip-network.pdf}


\begin{list2}
\item Network traffic is sent using network protocols with fields, values and destinations
\item Can we inspect this and decide to pass or block traffic, certainly! Does it help, sure!
\end{list2}


\slide{Protocols: OSI and Internet models}

\hlkimage{9cm,angle=90}{images/compare-osi-ip.pdf}

Protocols are rules, like how to say; Hello!
Don't worry too much yet, this is an overview!


\slide{Recommended network technologies to learn}

Small set of recommended networking technologies to learn

Networking: Basic Protocols from the Internet Protocols suite -- TCP/IP
\begin{list2}
\item Network Layer: Ethernet, Address Resolution Protocol (ARP), IPv4 and ICMPv4
\item Transport Layer: Transmission Control Protocol (TCP) and User Datagram Protocol (UDP)
\item Common upper layer: Dynamic Host Configuration Protocol (DHCP), Domain Name System (DNS),
Hypertext Transfer Protocol (HTTP)\\
\vskip 5mm
\item Later add the encrypted/secure versions like Hypertext Transfer Protocol Secure (HTTPS) which uses Transport Layer Security (TLS)
\item Later add Wi-Fi and IPv6/ICMPv6 with Neighbor Discovery Protocol (NDP)
\end{list2}

Enough to be able to go to other resources, and a connected path from ARP to TLS



\slide{Reversing and Attacking Network Protocols}

\hlkimage{4cm}{anp_cover-front-final.png}

A method with lots detail can be found in the book,\\
\emph{Attacking Network Protocols A Hacker's Guide to Capture, Analysis, and Exploitation}\\
by James Forshaw December 2017, 336 pp. ISBN-13: 9781593277505

\url{https://nostarch.com/networkprotocols}


\slide{Auditing Application Protocols}

\begin{list2}
\item Collect documentation
\item Identify Elements of Unknown Protocols
\item Use packet sniffers, tcpdump and Wireshark
\item Initiate the Connection Several Times
\item Replay traffic, can sometimes replay even encrypted traffic, see wireless WEP attacks
\end{list2}

Note: We investigate protocols, so we can see what is sent, so we can design \emph{payloads} which create problems for implementations - applications

\slide{Reverse Engineer Applications}

\begin{alltt}\footnotesize
  (gdb) disas main
  Dump of assembler code for function main:
     0x0000000000000580 <+0>:	lea    0x1ed(%rip),%rdi        # 0x774
     0x0000000000000587 <+7>:	sub    $0x8,%rsp
     0x000000000000058b <+11>:	mov    $0x7fff,%esi
     0x0000000000000590 <+16>:	xor    %eax,%eax
     0x0000000000000592 <+18>:	callq  0x560 <printf@plt>
     0x0000000000000597 <+23>:	lea    0x1ed(%rip),%rdi        # 0x78b
     0x000000000000059e <+30>:	mov    $0xffff8000,%esi
     0x00000000000005a3 <+35>:	xor    %eax,%eax
     0x00000000000005a5 <+37>:	callq  0x560 <printf@plt>
     0x00000000000005aa <+42>:	xor    %eax,%eax
     0x00000000000005ac <+44>:	add    $0x8,%rsp
     0x00000000000005b0 <+48>:	retq
  End of assembler dump.
\end{alltt}

\begin{list2}
\item It is possible to debug, disassemble and reverse engineer applications
\item Calling socket functions, seeing structs, data types etc.
\item Examine strings: HTTP, FTP, SMTP etc. all uses semi-english words GET, EHLO, PASS
\end{list2}


\slide{Special values}

\begin{list2}
\item Examine special values
\item What are the defined/used values
\item What happens if this is changed? Do they cover values outside of the used ranges? Case/switch constructs
\item Use trace functions in the operating system, can capture, analyze and replay sometimes
\end{list2}


\slide{Buffer Overflow when receiving}

\begin{list2}
\item When you see data enter the application, identify functions
\item Consider if they use dangerous functions, strcpy and friends
\item How much space is available, allocated etc.
\item Basic stuff and similar across applications
\vskip 2cm
\item Repeat everything we learned about string processing, integeroverflows/underflows etc. Just from the network
\item Often trying to abuse will lead to denial of service
\vskip 1cm
\item If some rock solid service starts bouncing down and up, maybe look into traffic received.
\item This is what honeypots also do
\end{list2}

\slide{Binary Protocols}

\begin{list2}
\item Some protocols use binary formats
\item Example DNS, which is a complex protocol
\item When parsing DNS use standard libraries!
\item When attacking DNS applications, use standard libraries! \smiley
\item DNS is just an example, new protocols may not be implemented in easy to use tools\\ -- but someone might have analyzed it or parts already!
\end{list2}

\slide{ASN.1 problems}

\begin{list2}
\item Abstract format designed for representing objects in a machine independent format
\item Used for various technologies in use on the internet:
\item Certificates and key encoding
\item Simple Network Management Protocol (SNMP)
\item ISAKMP part of IPsec
\item Lightweight Directory Access Protocol (LDAP)
\end{list2}

\slide{Linux Kernel ASN.1}

\begin{list2}
\item CVE-2016-0758 Integer overflow in lib/asn1\_decoder.c in the Linux kernel before 4.6 allows local users to gain privileges via crafted ASN.1 data.\\
\url{https://cve.mitre.org/cgi-bin/cvename.cgi?name=CVE-2016-0758}
\item Linux kernel have about 5 ASN.1 parsers\\
\url{https://www.x41-dsec.de/de/lab/blog/kernel_userspace/}
\end{list2}

\slide{ASN.1 parsing CVEs -- 2003 to 2019}

\begin{list2}
\item CVE-2003-0818 Microsoft MSASN1.DLL (MS04-007), CVSS 9.8:
      BER length integer overflows, RCE via Kerberos, NTLM, IPsec, TLS\\
\url{https://docs.microsoft.com/en-us/security-updates/securitybulletins/2004/ms04-007}
\item CVE-2016-2053 Linux kernel, CVSS 4.7: panic in \verb+asn1_ber_decoder+\\
\url{https://github.com/torvalds/linux/commit/0d62e9dd6da45bbf0f33a8617afc5fe774c8f45f}
\item CVE-2016-2108 OpenSSL, CVSS 9.8: ``negative zero'',
      memory corruption via crafted ANY field\\
\url{https://nvd.nist.gov/vuln/detail/CVE-2016-2108}
\item CVE-2016-5080 Objective Systems ASN1C (Cisco ASR 5000), CVSS 9.8:
      integer overflow in compiler-generated code\\
      \url{https://www.cisco.com/c/en/us/support/docs/csa/cisco-sa-20160721-asn1c.html}
\item CVE-2019-9162 Linux kernel, CVSS 7.8: missing ASN.1 length checks
      in SNMP NAT helper, OOB read/write\\
      \url{https://bugs.chromium.org/p/project-zero/issues/detail?id=1776}
\end{list2}

Need a rabbit hole, check out ASN.1 parsing, parsing anything is hard!

\slide{ASN.1 parsing CVEs -- 2021 to 2025}

\begin{list2}
\item CVE-2021-43527 Mozilla NSS ``BigSig'', CVSS 9.8: heap overflow on
      oversized DER-encoded signature\\
      \url{https://nvd.nist.gov/vuln/detail/CVE-2021-43527}
\item CVE-2022-44640 Heimdal Kerberos KDC, CVSS 9.8: invalid free on
      crafted DER CHOICE, generated by its ASN.1 compiler
      \url{https://github.com/heimdal/heimdal/security/advisories/GHSA-88pm-hfmq-7vv4}
\item CVE-2023-0286 OpenSSL, CVSS 7.4: X.400 address type confusion
      in GeneralName, attacker pointers to \verb+memcmp+\\
      \url{https://www.openssl.org/news/secadv/20230207.txt}
\item CVE-2024-12133 libtasn1, CVSS 5.3: many SEQUENCE OF/SET OF
      elements, quadratic decoding, DoS\\
      \url{https://gitlab.com/gnutls/libtasn1/-/issues/52}
\item CVE-2025-0343 Swift ASN.1, CVSS 7.5: primitive vs.\ constructed
      form not enforced, crash on crafted BER/DER\\
      \url{https://nvd.nist.gov/vuln/detail/CVE-2025-0343}
\item CVE-2025-13151 libtasn1, CVSS 7.5 (upstream: Low): one-byte
      stack overflow in \verb+asn1_expand_octet_string+\\
      \url{https://gitlab.com/gnutls/libtasn1/-/issues/55}
\end{list2}

\begin{list1}
\item Same bug class for 20+ years: trusting length, type and form fields
\end{list1}


\slide{Type Length Value TLVs}

\begin{quote}
  TLV sequences are easily searched using generalized parsing functions;
New message elements which are received at an older node can be safely skipped and the rest of the message can be parsed. This is similar to the way that unknown XML tags can be safely skipped;
TLV elements can be placed in any order inside the message body;
TLV elements are typically used in a binary format which makes parsing faster and the data smaller than in comparable text based protocols.
\end{quote}
Source: \url{https://en.wikipedia.org/wiki/Type-length-value}

\begin{list2}
\item Type Length Value is an encoding used in data communication
\item For example in Link Layer Discovery Protocol (LLDP), SSH, RADIUS
\end{list2}



\slide{Network Authentication}

\begin{quote}{\bf
  IPMI Authentication Bypass via Cipher 0}\\
  Dan Farmer identified a serious failing of the IPMI 2.0 specification, namely that cipher type 0, an indicator that the client wants to use clear-text authentication, actually allows access with any password. Cipher 0 issues were identified in HP, Dell, and Supermicro BMCs, with the issue likely encompassing all IPMI 2.0 implementations. It is easy to identify systems that have cipher 0 enabled using the \verb+ipmi_cipher_zero+ module in the Metasploit Framework.
\end{quote}

\begin{list2}
\item Sometimes people add network functionality to existing applications\\
-- and do this badly
\item We have seen applications like IPMI and others
\item Is this a network problem, server problem, hardware problem? -- {\bf discuss}
\end{list2}

Source: \url{https://blog.rapid7.com/2013/07/02/a-penetration-testers-guide-to-ipmi/}

\slide{Even securiy protocols and software have problems}

\begin{list2}
\item ISAKMP example
\item IKE(v1) has been critized as being overly complex
\item Needed bake-off sessions where vendors meet and tried negotiating
\item Searching for CVE ISAKMP show multiple vulnerabilities in various implementations, including firewalls and tcpdump
\item Source: AoSSA chapter 16: Network Application Protocols\\
\emph{The Art of Software Security Assessment Identifying and Preventing
Software Vulnerabilities}
Mark Dowd, John McDonald, Justin Schuh ISBN: 9780321444424
\end{list2}

\breakdance

\slide{Session 1: Exercise: *sshole}

\slide{password: AAAAAAAAAAAAAAAA...AAAAAAAAAAAAA}

\begin{quote}{\bf\large
*sshole}

so openssh on musl linux (alpine, openwrt, buildroot, yocto, and a lot of the docker images you've pulled) has had a little problem since 2016.

accounts that should never take a password at all (root, your key-only deploy user, ldap users, the default user on alpine's aws and azure images) take any password longer than 256 characters.

yes, any password. yes. i know. yes.
\end{quote}
Source: \url{https://sshhole.com/}!

Lets look into this including CVE-2016-6210 and CVE-2019-5021

\slide{Release notes: OpenSSH 10.6p2 (2026-11-05)}

New version
\begin{alltt}\small
Portable OpenSSH 10.6p2 was released on 2026-11-05. It will be available
from the mirrors listed at https://www.openssh.com/.
...
Security
========
 * sshd(8): \textbf{avoid password authentication bypass} that is \textbf{exploitable}
   through sshd(8) that affects some Linux systems when built and run
   with \textbf{PAM support disabled}. We consider the root cause of this bug
   to be \textbf{outside sshd}, but it is \textbf{exploitable via sshd} and \textbf{can be
   avoided there}.

   This bug was reported by Özgün Kültekin of Wakefield
   (wakefield.co)
\end{alltt}
Source: \url{https://www.openssh.org/releasenotes.html} bold by me

\slide{OpenSSH receives large number of security bug reports}

\begin{alltt}\small
Recently the OpenSSH team have received a large number of security
bug reports, many of which are findings from AI models or made with
AI assistance. While many AI reports are determined not to have
security impact when considered in the context of a realistic
threat model, we very much welcome these reports, especially when
combined with human triage, analysis, test-cases and particularly
when accompanied by proposed fixes.
\end{alltt}
Source: \url{https://www.openssh.org/releasenotes.html}


\slide{OpenSSH AI changes disclosure policies!}

\begin{alltt}\small
** We have seen a number of cases where a security bug identified
** by AI tools is subsequently independently discovered by a
** different researcher. This suggests that adversaries who do not
** report bugs to OSS projects are likely to be able to discover
** these bugs too. Given this, the OpenSSH team will, for now, be
** making more frequent releases to get bugfixes into users' hands
** more quickly rather than batching them until the next planned
** release.
\end{alltt}
Source: \url{https://www.openssh.org/releasenotes.html}

Saw something similar about cURL!

\slide{Lets try it!}

\begin{list2}
\item Read up on the bugs -- multiple, note how it would have been hard predicting the interactions!
\item Lets find the container images -- run them in your K8s clusters
\item Maybe it needs some configuration?
\vskip 5mm
\item When the attack works, retry with updated software/images -- and same configuration
\end{list2}

Bonus question: would any of our previous Defense in Depth strategies work? Running container as non-root?

Would VictoriaLogs or similar work? What would that provide us?

\breakdance


\slide{Sessions 2+3: Teaching and exercises in Networks and DDoS}


\slide{Networking with some pentesting}

{\bf What is pentest}

\begin{quote}
A penetration test, informally pen test, is an attack on a computer system that looks for security weaknesses, potentially gaining access to the computer's features and data.[1][2]
\end{quote}
Source: quote from \link{https://en.wikipedia.org/wiki/Penetration_test}

\begin{list1}
\item Penetration testing is a simulation, with good intentions
\item People around the world constantly \emph{test your defenses}
\item Often better to test at planned times
\item How to create DDoS simulations, tools and process
\item I use and recommend Kali Linux as the base for this
\end{list1}


\slide{Networks today}
\hlkimage{13cm}{overview-routing-customer-2015.pdf}

\begin{list1}
\item Conclusion: networks are complex
\item Many places where you can block, filter, restrict, limit, avoid, process, change, observe
\end{list1}

\slide{Book: Applied Network Security Monitoring (ANSM)}

\hlkimage{5cm}{ansm-book.png}

\emph{Applied Network Security Monitoring: Collection, Detection, and Analysis}
1st Edition

Chris Sanders, Jason Smith
eBook ISBN: 9780124172166
Paperback ISBN: 9780124172081 496 pp.
Imprint: Syngress, December 2013

{\footnotesize\link{https://www.elsevier.com/books/applied-network-security-monitoring/unknown/978-0-12-417208-1}}

\slide{Book: Practical Packet Analysis (PPA)}
\hlkimage{6cm}{PracticalPacketAnalysis3E_cover.png}

\emph{Practical Packet Analysis,
Using Wireshark to Solve Real-World Network Problems}
by Chris Sanders, 3rd Edition
April 2017, 368 pp.
ISBN-13:
978-1-59327-802-1

\link{https://nostarch.com/packetanalysis3}


\slide{Internet is Open Standards!}

\begin{quote}
We reject kings, presidents, and voting.\\
We believe in rough consensus and running code.\\
-- The IETF credo Dave Clark, 1992.
\end{quote}

\begin{list1}
\item Request for comments - RFC - er en serie af dokumenter
\item RFC, BCP, FYI, informational\\
de første stammer tilbage fra 1969
\item Ændres ikke, men får status Obsoleted når der udkommer en nyere
  version af en standard
\item Standards track:\\
Proposed Standard $\rightarrow$ Draft Standard $\rightarrow$ Standard
\item  Åbne standarder = åbenhed, ikke garanti for sikkerhed
\end{list1}




\slide{IPv4 packets -- header - RFC-791}

\begin{alltt}
\small
    0                   1                   2                   3
    0 1 2 3 4 5 6 7 8 9 0 1 2 3 4 5 6 7 8 9 0 1 2 3 4 5 6 7 8 9 0 1
   +-+-+-+-+-+-+-+-+-+-+-+-+-+-+-+-+-+-+-+-+-+-+-+-+-+-+-+-+-+-+-+-+
   |Version|  IHL  |Type of Service|          Total Length         |
   +-+-+-+-+-+-+-+-+-+-+-+-+-+-+-+-+-+-+-+-+-+-+-+-+-+-+-+-+-+-+-+-+
   |         Identification        |Flags|      Fragment Offset    |
   +-+-+-+-+-+-+-+-+-+-+-+-+-+-+-+-+-+-+-+-+-+-+-+-+-+-+-+-+-+-+-+-+
   |  Time to Live |    Protocol   |         Header Checksum       |
   +-+-+-+-+-+-+-+-+-+-+-+-+-+-+-+-+-+-+-+-+-+-+-+-+-+-+-+-+-+-+-+-+
   |                       Source Address                          |
   +-+-+-+-+-+-+-+-+-+-+-+-+-+-+-+-+-+-+-+-+-+-+-+-+-+-+-+-+-+-+-+-+
   |                    Destination Address                        |
   +-+-+-+-+-+-+-+-+-+-+-+-+-+-+-+-+-+-+-+-+-+-+-+-+-+-+-+-+-+-+-+-+
   |                    Options                    |    Padding    |
   +-+-+-+-+-+-+-+-+-+-+-+-+-+-+-+-+-+-+-+-+-+-+-+-+-+-+-+-+-+-+-+-+

                    Example Internet Datagram Header
\end{alltt}


\slide{The basic tools for looking at packets}

Knowledge and insight
\begin{list2}
\item Networks have end-points and conversations on multiple layers
\item Name resolution includes low level MAC addresses, and IP -- names
\item Tcpdump format, built-in to many network devices
\item Wireshark is advanced, try right-clicking different places
\item Remote packet dumps, like \verb+tcpdump –i eth0 –w packets.pcap+
\item Story: tcpdump was originally written in 1988 by Van Jacobson, Sally Floyd, Vern Paxson and Steven McCanne who were, at the time, working in the Lawrence Berkeley Laboratory Network Research Group\\
 \link{https://en.wikipedia.org/wiki/Tcpdump}
\end{list2}

\vskip 5mm

You can learn a lot by just listening!


\slide{Wireshark - graphical network sniffer}

\hlkimage{13cm}{images/wireshark-http.png}

\centerline{Capture - Options, select a network interface}
\centerline{\link{http://www.wireshark.org}}



\slide{DNS problems}

\begin{quote}
The Domain Name System (DNS) [32][33] provides for a distributed database mapping host names to IP
addresses. An intruder who interferes with the proper operation of the DNS can mount a variety of
attacks, including denial of service and password collection. There are a number of vulnerabilities.
\end{quote}
Source: \emph{Security Problems in the TCP/IP Protocol Suite}, S.M. Bellovin 1989\\
\link{https://www.cs.columbia.edu/~smb/papers/ipext.pdf}

\begin{list1}
\item We have a lot of the same problems in DNS today
\item Plus some more caused by middle-boxes, NAT, DNS size, DNS inspection
\begin{list2}
\item DNS must allow both UDP and TCP port 53
\item Your DNS servers must have updated software, which use updated protocol behaviour \\
see DNS flag day \link{https://dnsflagday.net/}
\item DNS is unencrypted
\end{list2}
\end{list1}

\slide{Recommended Software -- Unbound}

\begin{quote}
Unbound is a validating, recursive, caching DNS resolver. It is designed to be fast and lean and incorporates modern features based on open standards.

To help increase online privacy, Unbound supports DNS-over-TLS which allows clients to encrypt their communication. In addition, it supports various modern standards that limit the amount of data exchanged with authoritative servers.
\end{quote}

\link{https://www.nlnetlabs.nl/projects/unbound/about/}

My preferred local DNS server


\slide{DNS problems}

\begin{list2}
\item Failure to Deal with Invalid Label Lengths
\item Insufficient Destination Lengths Check
\item Insufficient Source Length Checks
\item Pointer Values Not Verified In Packet
\item Special Pointer Values
\item Length Variables
\item These are from the book: AoSSA chapter 16: Network Application Protocols\\
\emph{The Art of Software Security Assessment Identifying and Preventing
Software Vulnerabilities}
Mark Dowd, John McDonald, Justin Schuh ISBN: 9780321444424
\end{list2}

\vskip 1cm
\centerline{Labels and pointers within packets save bytes, but make it more complex!}


\slide{Firewalls and related issues}

\begin{quote}
In computing, a firewall is a network security system that monitors and controls incoming and outgoing network traffic based on predetermined security rules.[1] A firewall typically establishes a barrier between a trusted internal network and untrusted external network, such as the Internet.[2]
\end{quote} Source: Wikipedia

\begin{list2}
\item \link{https://en.wikipedia.org/wiki/Firewall_(computing)}
\item \link{http://www.wilyhacker.com/} Cheswick chapter 2 PDF
\emph{A Security Review of Protocols:
Lower Layers}
\item Network layer, packet filters, application level, stateless, stateful
\end{list2}

Firewalls are by design a choke point, natural place \\
to do network security monitoring!



\slide{Back in the day: Firewalls and Internet Security, 1994}

\hlkimage{4cm}{images/cheswick-cover2e.jpg}
\begin{list2}
\item \emph{Firewalls and Internet Security: Repelling the Wily Hacker} , Second Edition 2003, William R. Cheswick, Steven M. Bellovin, and Aviel D. Rubin, {\bf 2003} \link{http://www.wilyhacker.com/}
\item The full PDF—and the full LaTeX source of the book.
Creative Commons Attribution-NonCommercial-NoDerivatives 4.0 International License.
\item How to configure firewalls often boil down to, should we allow protocol X
\item If we allow certain protocols through a firewall, we are asking for trouble
\end{list2}



\slide{Definition of firewalls -- multiple definitions exist}

\begin{quote}
We define a firewall as {\bf a collection of components} placed between two networks that collectively have
the following properties:
\begin{list2}
\item All traffic from inside to outside, and vice-versa, must pass through the firewall.
\item Only authorized traffic, as defined by the {\bf local security policy}, will be allowed to pass.
\item The firewall itself is immune to penetration.
\end{list2}
We should note that these are design goals; a failure in one aspect does not mean that the collection
is not a firewall, simply that it is not a very good one.
\end{quote}

We will consider this a firewall, but we know today that {\bf both inside and outside at meaningless}, since we have {\bf multiple networks inside, we have partner network connections etc.}

Source: \citetitle{Cheswick94} by \citeauthor{Cheswick94} \citeyear{Cheswick94}

\slide{Definition of firewalls -- Wikipedia}

Another short definition that encapsulates this is found on Wikipedia, and may suffice in many situations. Again there will typically be multiple networks, zones or areas of the networks with varying degrees of trust.
\begin{quote}
In computing, a firewall is a {\bf network security system that monitors and controls incoming and outgoing network traffic based on predetermined security rules}.[1] A firewall typically establishes a barrier between a trusted network and an untrusted network, such as the Internet.[2]
\end{quote}
Source: Wikipedia
 \link{https://en.wikipedia.org/wiki/Firewall_(computing)}

{\bf TL;DR Not necessarily a single device}



\slide{A firewall -- in the vendor eyes}

% Single line in firewall - single line out
% With text device under test, as caption
\hlkimage{19cm}{Firewall-vendor.pdf}


\begin{quote}
"{\bf Can your firewall flex in the face of change?}\\
Does it harmonize your network, workload, and application security? Does it protect apps and employees in your hybrid or multicloud environment? Make sure you're covered."
\end{quote}
Source: not shown to protect the audience from further marketing speak
% Source is Cisco.com




\slide{A firewall -- in the enterprise mindset }

% Single line in firewall,
% cloud with at-sign Internet on the left,\\
% cloud with at-sign LAN on the right
\hlkimage{19cm}{Firewall-enterprise.pdf}

\begin{list2}
\item Even though some vendors suggest they can do everything in a single box, I don't believe them!
\item Truth -- yes, we can do almost anything in software
\item Realization {\bf Your infrastructure is based on multiple components and or devices}
\end{list2}


\slide{Defense in depth}

%\hlkimage{10cm}{Bartizan.png}
\hlkimage{15cm}{medieval-clipart-5}
\centerline{Picture originally from: \url{http://karenswhimsy.com/public-domain-images}}

\slide{Firewalls}

\hlkimage{8cm}{opal-luci-firewall-easy.png}

\begin{list2}
\item Firewalls are a great project to look into, there are so many \smiley
\item I recommend looking into open source firewalls like OPNsense, UFW and OpenBSD and many routers also include firewalls.
 You can often find images from commercial vendors, that you can try out.
\item Picture is from GL-Inet router Opal with OpenWRT/LuCI which is router software for smaller routers
\end{list2}


\slide{What is a packet filter}

We may want to consider the difference between different types of firewalls/devices:
\begin{list2}
\item Network layer, we often call them packet filters, stateless
\item Application level, we often call them, stateful filtering and gateways
\item Firewalls are by design a choke point, natural place \\
to do network security monitoring!
%\item Older but still interesting Cheswick chapter 2 PDF
%\emph{A Security Review of Protocols: Lower Layers}\\
%\link{http://www.wilyhacker.com/}
\end{list2}

They are all firewalls -- or firewall devices!



\slide{Firewall OPNsense GUI based and easy to install}

\hlkimage{12cm}{images/screenshots_OPNsense-1024x518.png}

\begin{list2}
\item OPNsense \link{https://opnsense.org/}
\item Firewall built on FreeBSD with web interface
\item Originally thoughts from m0n0wall and later \link{https://www.pfsense.org/}\\
\item Danish companies have been using these for many years now
\end{list2}



\slide{Linux Uncomplicated Firewall (UFW)}

Debian / Ubuntu:
\begin{alltt}\scriptsize
root@debian01:~# apt install ufw
...
root@debian01:~# ufw allow 22/tcp
Rules updated
Rules updated (v6)
root@debian01:~# ufw enable
Command may disrupt existing ssh connections. Proceed with operation (y|n)? y
Firewall is active and enabled on system startup
root@debian01:~# ufw status numbered
Status: active

     To                         Action      From
     --                         ------      ----
[ 1] 22/tcp                     ALLOW IN    Anywhere
[ 2] 22/tcp (v6)                ALLOW IN    Anywhere (v6)
\end{alltt}

\begin{list2}
\item Extremely easy to use -- I recommend and use the (Uncomplicated Firewall) UFW
\end{list2}





\slide{Sample rules from OpenBSD PF}

\begin{alltt}\tiny
# hosts and networks
router="217.157.20.129"
webserver="217.157.20.131"
homenet="{ 192.168.1.0/24, 1.2.3.4/24 }"
wlan="10.0.42.0/24"
wireless=wi0
set skip lo0
# things not used
spoofed="{ 127.0.0.0/8, 172.16.0.0/12, 10.0.0.0/16, 255.255.255.255/32 }"
{\bf
block in all # default block anything}
# egress and ingress filtering - disallow spoofing, and drop spoofed
block in quick from $spoofed to any
block out quick from any to $spoofed

pass in on $wireless proto tcp from \{ $wlan $homenet \} to any port = 22
pass in on $wireless proto tcp from any to $webserver port = 80

pass out
\end{alltt}



\slide{Linux TCP SACK PANIC  - CVE-2019-11477 et al}

Kernel vulnerabilities, CVE-2019-11477, CVE-2019-11478 and CVE-2019-11479

\begin{quote}\footnotesize{\bf
Executive Summary}\\
Three related flaws were found in the Linux kernel’s handling of TCP networking.  The most severe vulnerability could allow a remote attacker to trigger a kernel panic in systems running the affected software and, as a result, impact the system’s availability.

The issues have been assigned multiple CVEs: CVE-2019-11477 is considered an Important severity, whereas CVE-2019-11478 and CVE-2019-11479 are considered a Moderate severity.

The first two are related to the Selective Acknowledgement (SACK) packets combined with Maximum Segment Size (MSS), the third solely with the Maximum Segment Size (MSS).

These issues are corrected either through applying mitigations or kernel patches.  Mitigation details and links to RHSA advsories can be found on the RESOLVE tab of this article.
\end{quote}

Source: {\footnotesize\url{https://access.redhat.com/security/vulnerabilities/tcpsack}}


\slide{Linux TCP IPv4 options UAF - CVE-2026-98432}

Kernel vulnerability CVE-2026-98432, present since Linux 4.4 (2016)

\begin{quote}\footnotesize{\bf
Summary}\\
A race condition in the Linux kernel's handling of new IPv4 TCP connections can lead to a use-after-free. IP options stored with a pending connection request are handed to the new socket without being copied. If that socket is destroyed while a retransmitted SYN is still being answered, freed memory is used. IPv6 is not affected.

An unauthenticated remote attacker sends SYNs carrying IPv4 options (Record Route or Timestamp) to a listening port and retries until the race is won. A lost race only fails that one connection attempt, so the attacker can keep trying.

Anthropic demonstrated remote root against servers with a Multipath TCP listener. Go 1.24 and later programs open MPTCP listeners by default.

Fixed in 5.10.272, 5.15.223, 6.1.190, 6.6.159, 6.12.113, 6.18.56 and 7.2.10. As a mitigation, Red Hat suggests filtering incoming IPv4 TCP packets that carry header options.
\end{quote}

Source: {\footnotesize\url{https://access.redhat.com/security/vulnerabilities/RHSB-2026-016}}

\slide{Rooting devices from far away}

\begin{quote}\small{\bf
\# Root from kilometers away: RCE in Ubiquiti airMAX antennas}\\
Researchers at Faraday Security (Gastón Aznarez, Federico Kirschbaum, Dan Borgogno) found two unauthenticated remote code execution vulnerabilities in Ubiquiti's outdoor wireless equipment: airMAX and airFiber antennas and UniFi building bridges. These are the devices mounted on rooftops and poles to carry internet, cameras and building links over long distances.

**CVEs** (both CVSS 3.1 score 8.8, `AV:A/AC:L/PR:N/UI:N/S:U/C:H/I:H/A:H`):

- **CVE-2026-21639**: airMAX AC, airMAX M, airFiber AF60-XG, airFiber AF60\\
- **CVE-2026-21638**: UBB-XG, UDB-Pro/UDB-Pro-Sector, UBB

\end{quote}
Source: \url{https://faradaysec.com/update-your-ubiquiti-equipment-what-we-found-in-airmax-antennas/}

\slide{Why?!}
\begin{quote}\small{\bf
\# What went wrong}

- The proprietary airMAX protocol is carried inside vendor-specific fields of ordinary Wi-Fi management frames.
- Parts of that data looked encrypted, but were only obfuscated.
- The flaw is a command injection (CWE-77): input from the radio frame reached system command execution without proper sanitisation.
- Several affected models ran \textbf{internal software around 17 years old, unsupported for close to a decade.}
\end{quote}

\begin{list2}
\item Faraday blog post: \url{https://faradaysec.com/update-your-ubiquiti-equipment-what-we-found-in-airmax-antennas/}
\item Ubiquiti Security Advisory Bulletin 061 (CVE-2026-21639):\\
\url{https://community.ui.com/releases/Security-Advisory-Bulletin-061-061/1e4fe5f8-29c7-4a7d-a518-01b1537983ba}
\item Ubiquiti Security Advisory Bulletin 060 (CVE-2026-21638):\\
 \url{https://community.ui.com/releases/Security-Advisory-Bulletin-060-060/cde18da7-2bc4-41bb-a9cc-48a4a4c479c1}
\item The talk "Root From Kilometers Away" was presented at Black Hat USA 2026 and DEF CON 34.
\end{list2}


\slide{Bottlenecks exist, but where}


\hlkimage{12cm}{overview-routing-customer-2015.pdf}

\begin{list2}
\item Lower layer attacks Transport Layer Attacks TCP SYN flood -- packet based
\item Higher layer attacks like Slowloris and web attacks -- keep sessions running
\item Protect everything without loosing functionality or creating administrative nightmare
\end{list2}

\slide{Availability and Network flooding attacks}

\begin{list2}
\item {\bf SYN flood} is the most basic and very common on the internet towards 80/tcp and 443/tcp
\item {\bf ICMP and UDP flooding} are the next popular targets -- more similar ones exist
\item Special packets and protocols -- anything that can create \emph{load on systems} work
\item All of them try to use up some resources
\begin{list2}
\item {\bf Memory space} in specific sections of the {\bf kernel, TCP state, firewalls state, number of concurrent sessions/connections}
\item {\bf Interrupt processing} of packets - packets per second (pps)
\item {\bf CPU processing} in firewalls, pps
\item CPU processing in server software
\item {\bf Bandwidth} - megabits per second (mbps)
\item Typically source is spoofed or amplification attacks abusing devices on the Internet
\end{list2}
\end{list2}


\slide{Packet processing in firewalls -- detailed view}

% Picture with pipeline from Juniper SRX
\hlkimage{23cm}{srx-firewall-flow-based-processing.png}
\emph{Traffic Processing on SRX Series Devices Overview}\\ {\scriptsize
\link{https://www.juniper.net/documentation/us/en/software/junos/flow-packet-processing/topics/topic-map/security-srx-devices-processing-overview.html}}



\slide{Simulating DDoS packets}

A minimal introduction workshop teaching people how to produce DDoS simulation traffic - usefull for testing their own infrastructures.

We will have a server connected on a switch with multiple 1Gbit port for attackers. Attackers will be connected through 1Gbit ports using USB Ethernet - we have loaners.

Work together to produce enough to take down this server!

WHILE attack is ongoing there will be both the possibility to monitor traffic, monitor port, and decide on changes to prevent the attacks from working.

\slide{Common DDoS attack types}

We will work through common attack types, like:

\begin{list2}
\item TCP SYN flooding
\item TCP other flooding
\item UDP flooding NTP, etc.
\item ICMP flooding
\item Misc - stranger attacks and illegal combinations of flags etc.
\end{list2}

then we will discuss which changes to environment could be implemented.

You will go away from this with tools for producing packets, hping3 and some configurations for protecting - PF rules, switch rules, server firewall rules.



\slide{Scanning and Attacking -- \bf Pressure Points and Scope}

% Drawing showing "network"

\hlkimage{21cm}{generic-network-pressure-points.pdf}

\begin{list2}
\item In scope for me is everything that could adversely affect the network
\item Typical scope IPv4: Link network /28 or /26 and a hosting network /24 or even /22
\item Typical scope IPv6: Link network /64 (bad) or /127 (RFC9099) and a hosting network /48 with subnets
%\item Domain Name Service outside of the network is ofc also attack vector, ignored today
\end{list2}



\slide{Prepare for the testing}


\hlkimage{10cm}{trinity-nmapscreen-hd-cropscale-418x250.jpg}

\begin{list2}
\item Portscan the whole linknet and hosting range in IPv4
\item Ask about IPv6 ranges in use, specific subnets and IPv6 addresses\\
We can guess from traceroute, Nmap test first 100 addresses in each subnet etc. but easier to ask
\item Portscan the whole linknet - identify provider devices, hosting network devices, type of device router/firewall
\item Also Thank you Fyodor (Gordon Lyon) and contributors for Nmap!
\end{list2}


\slide{Detailed Scope and Plan}

Select a few targets for monitoring and attacks, from the port scans

Best case would be to have:
\begin{list2}
\item Ping ICMP allowed to a provider router and hosting firewall -- is the connection alive
\item TCP port with service checks, HTTP being attacked and one not being attacked
\item UDP port with service check, DNS is a favourite -- ask for localhost/127.0.0.1
\item Put monitoring on these, a week before testing is nice
\item Agree on a day or night for testing, inform participants and system owners
\end{list2}


\slide{hping3 packet generator}

\begin{alltt}\footnotesize
usage: hping3 host [options]
  -i  --interval  wait (uX for X microseconds, for example -i u1000)
      --fast      alias for -i u10000 (10 packets for second)
      --faster    alias for -i u1000 (100 packets for second)
      --flood	   sent packets as fast as possible. Don't show replies.
...
hping3 is fully scriptable using the TCL language, and packets
can be received and sent via a binary or string representation
describing the packets.
\end{alltt}

\begin{list2}
\item Hping3 packet generator is a very flexible tool to produce simulated DDoS traffic with specific charateristics
\item Home page: \link{http://www.hping.org/hping3.html}
\item Source repository \link{https://github.com/antirez/hping}
\end{list2}

\centerline{My primary DDoS testing tool, easy to get specific rate pps}


\slide{t50 packet generator}


\begin{alltt}\footnotesize
root@cornerstone03:~# t50 -?
T50 Experimental Mixed Packet Injector Tool 5.4.1
Originally created by Nelson Brito <nbrito@sekure.org>
Maintained by Fernando Mercês <fernando@mentebinaria.com.br>

Usage: T50 <host> [/CIDR] [options]

Common Options:
    --threshold NUM        Threshold of packets to send     (default 1000)
    --flood                This option supersedes the 'threshold'
...
6. Running T50 with '--protocol T50' option, sends ALL protocols sequentially.
root@cornerstone03:~# t50 -? | wc -l
264
\end{alltt}

\begin{list2}
\item T50 packet generator, another high speed packet generator
can easily overload most firewalls by producing a randomized traffic with multiple protocols like IPsec, GRE, MIX \\
home page: \link{http://t50.sourceforge.net/resources.html}
\end{list2}

\centerline{Fast and breaks most firewalls when flooding, easy 800k pps/400Mbps}


\slide{Enter MoonGen and Dedicated Hardware}

\hlkimage{6cm}{dell-precision-3240-compact.jpg}

\begin{list2}
\item Modern computer with modern CPU and PCIe x8 or better\\
I currently use a few cheap Dell devices Precision 3640 Tower / Precision 3240 Compact (not a recommendation)
\item Supported card - I am using the old Intel 82599 based 10Gbit cards
\item DPDK and software -- I use MoonGen \link{https://github.com/emmericp/MoonGen}
\item A huge thanks to Paul Emmerich emmericp for programming and publishing his works!
\item Maybe the easiest way to use DPDK currently -- "Craft all packets in user-controlled Lua scripts"
\end{list2}


\slide{Penguingping packet generator}

\hlkimage{13cm}{penguinping-02-peak.png}

\begin{list2}
\item PenguinPing packet generator, my high speed packet generator
home page: \link{https://penguinping.org}
\item First versions are only about 230 lines of Lua code and implement basic command line to replace hping3
\item Built on top of MoonGen/libmoon \link{https://github.com/emmericp/MoonGen}
\end{list2}

\centerline{Extremely fast and allows easy customization}


\slide{Process: monitor, attack, break, repeat}

\begin{list2}
\item Pre-test: Monitoring setup - from multiple points
\item Pre-test: Perform full Nmap scan of network and ports
\item Start small, run with delays between packets
\item Turn up until it breaks, decrease delay - until using \verb+--flood+
\item Monitor speed of attack on your router interface pps/bandwidth
\item Give it maximum speed\\
 \verb+hping3 --flood -1+ and \verb+hping3 --flood -2+
\item Have a common chat with network operators/customer to talk about symptoms and things observed
\item Any information resulting from testing is good information
\end{list2}

\vskip 1cm
\centerline{Ohh we lost our VPN into the environment, ohh the fw console is dead}

\slide{Before testing: Smokeping}

\hlkimage{15cm}{smokeping-before-testing.png}

\centerline{Before DDoS testing  use Smokeping software}

\slide{Before testing: Pingdom}

\hlkimage{18cm}{forside-pingdom.png}

\centerline{Another external monitoring from Pingdom.com}

\slide{Performing the DDoS test }

So lets break this task down into:
\begin{list2}
\item Use Nmap to port scan the network
\item Run testing tool \verb+hping3+, \verb+t50+ or \verb+penguinping+
\item Go through all scenarios before making changes to environment
\item Expect things to break, investigate, repeat failed scenarios
\vskip 1cm
\item BTW we usually schedule this for night time! There WILL be interruptions
\vskip 1cm
\item How do you run the tools, since they all have so many options?\\
\verb+man nmap | enscript -o test.ps+ result in 54 pages,
\verb+hping3+ has about 80 options and \verb+t50+ has even more
\end{list2}

\slide{Running full port scan on network}

\begin{alltt}
\small
# export CUST_NET="192.0.2.0/24"
# nmap -p 1-65535 -A -oA full-scan $CUST_NET
\end{alltt}

Performs a full port scan of the network, all ports

Saves output in "all formats" normal, XML, and grepable formats

Goal is to enumerate the ports that are allowed through the network.

Note: This command is pretty harmless, if something dies, then it is\\
\emph{vulnerable to normal traffic} - and should be fixed!


\slide{Running Attacks with hping3}

\begin{alltt}\small
# export CUST_IP=192.0.2.1
# date;time hping3 -q -c 1000000  -i u60 -S -p 80  $CUST_IP
 \end{alltt}

\begin{alltt}\small
# date;time hping3 -q -c 1000000  -i u60 -S -p 80  $CUST_IP
Thu Jan 21 22:37:06 CET 2016
HPING 192.0.2.1 (eth0 192.0.2.1): S set, 40 headers + 0 data bytes

--- 192.0.2.1 hping statistic ---
1000000 packets transmitted, 999996 packets received, 1% packet loss
round-trip min/avg/max = 0.9/7.0/1005.5 ms

real	1m7.438s
user	0m1.200s
sys	0m5.444s
\end{alltt}

\vskip 1cm
\centerline{Dont forget to do a killall hping3 when done \smiley }

\slide{Recommendations During Test}

\begin{list1}
\item Run each test for at least 5 minutes, or even 15 minutes\\
Some attacks require some build-up before resource run out
\item Take note of any change in response, higher latency, lost probes
\item If you see a change, then re-test using the same parameters, or a little less first
\item We want to know the approximate level where it breaks
\item If you want to change environment, then wait until all scenarios tested
\end{list1}

\slide{Comparable to real DDoS?}

Tools are simple and widely available but are they actually producing same result as high-powered and advanced criminal botnets. We can confirm that the attack delivered in this test is, in fact, producing the traffic patterns very close to criminal attacks in real-life scenarios.

\begin{list2}
\item We can also monitor logs when running a single test-case
\item Gain knowledge about supporting infrastructure
\item Can your syslog infrastructure handle 800.000 events in $<$ 1 hour?
\end{list2}


\slide{Running the tools}

A basic test would be:
\begin{list2}
\item TCP SYN flooding
\item TCP other flags, PUSH-ACK, RST, ACK, FIN
\item ICMP flooding
\item UDP flooding
\item Spoofed packets src=dst=target \smiley
\item Small fragments
\item Bad fragment offset
\item Bad checksum
\item Be creative
\item Mixed packets - like \verb+t50 --protocol T50+
\item Perhaps esoteric or unused protocols, GRE, IPSec
\end{list2}

\slide{Test-cases / Scenarios}

\begin{list1}
\item The minimal run contains at least these:
\begin{list2}
\item SYN flood: \verb+hping3 -q -c 1000000  -i u60 -S -p 80  $CUST_IP &+
\item SYN+ACK: \verb+hping3 -q -c 1000000  -i u60 -S -A -p 80  $CUST_IP &+
\item ICMP flood: \verb+hping3 -q -c --flood -1 $CUST_IP &+
\item UDP flood: \verb+hping3 -q -c --flood -1 $CUST_IP &+
\end{list2}
\item Vary the speed using the packet interval \verb+-i u60+ up/down
\item Use flooding with caution, runs max speeeeeeeeeeeed \smiley
\item TCP testing use a port which is allowed through the network, often 80/443
\item Focus on attacks which are hard to block, example TCP SYN must be allowed in
\item Also if you found devices like routers in front of environment\\
\verb+hping3 -q -c 1000000  -i u60 -S -p 22 $ROUTER_IP+\\
\verb+hping3 -q -c 1000000  -i u60 -S -p 179 $ROUTER_IP+
\end{list1}

\slide{Rocky Horror Picture Show - 1}

\hlkimage{20cm}{smokeping-1.png}

\centerline{Really does it break from 50.000 pps SYN attack?}

\slide{Rocky Horror Picture Show - 2}

\hlkimage{20cm}{smokeping-2.png}

\centerline{Oh no 500.000 pps UDP attacks work?}

\slide{Rocky Horror Picture Show - 3}

\centerline{Oh no spoofing attacks work?}

\hlkimage{20cm}{smokeping-3.png}


\exercise{ex:nmap-synscan}

\exercise{ex:nmap-pingsweep}

\exercise{ex:syn-flood}

Exercise booklet contains some bonus exercises, feel free to try them at home



\slide{Improvements seen after testing}

\begin{list1}
\item Turning off unneeded features - free up resources
\item Tuning sesions, max sessions src / dst
\item Tuning firewalls, max sessions in half-open state, enabling services
\item Tuning network, drop spoofed src from inside net \smiley
\item Tuning network, can follow logs, manage network during attacks
\item ...
\item And organisation has better understanding of DDoS challenges
\item Including vendors, firewall consultants, ISPs etc.
\end{list1}

\vskip 1cm
\centerline{After tuning of {\bf existing devices/network} improves results 10-100 times}

\slide{Conclusion}

\begin{list1}
\item You really should try testing
\item Investigate your existing devices\\
all of them, RTFM, upgrade firmware
\item Choose which devices does which\\
part - discard early to free resources\\
for later devices to dig deeper
\end{list1}

\vskip 2cm
\centerline{And dont forget that DDoS testing is as much a firedrill for the organisation}


%\slide{Stateless filtering Junos}
\slide{Stateless firewall filter throw stuff away}

\begin{alltt}\footnotesize
hlk@MX-CPH-02> show configuration firewall filter all | no-more
/* This is a static sample, perhaps better to use BGP flowspec and RTBH */
term edgeblocker \{
    from \{
        source-address \{
            84.180.xxx.173/32;
            87.245.xxx.171/32;
        \}
        destination-address \{
            91.102.91.16/28;
        \}
        protocol [ tcp udp icmp ];
    \}
    then \{
        count edge-block;
        discard;
    \}
\}
\end{alltt}
Hint: can also leave out protocol and then it will match all protocols

\slide{Stateless firewall filter limit protocols}

\begin{alltt}\footnotesize
term limit-icmp \{
    from \{
        protocol icmp;
    \}
    then \{
        policer ICMP-100M;
        accept;
    \}
\}
term limit-udp \{
    from \{
        protocol udp;
    \}
    then \{
        policer UDP-1000M;
        accept;
    \}
\}
\end{alltt}

Routers have extensive Class-of-Service (CoS) tools today

\slide{Strict filtering for some servers, still stateless!}

\begin{alltt}\footnotesize
term some-server-allow \{
    from \{
        destination-address \{
            109.238.xx.0/xx;
        \}
        protocol tcp;
        destination-port [ 80 443 ];
    \}
    then accept;
\}
term some-server-block-unneeded \{
    from \{
        destination-address \{
            109.238.xx.0/xx;
        \}
        protocol-except icmp;
    \}
    then discard;
\}
\end{alltt}

Wut - no UDP, yes UDP service is not used on these servers


\slide{ Firewalls - screens, IDS like features}

When you know regular traffic you can decide \emph{normal settings}:

\begin{alltt}\footnotesize
hlk@srx-kas-05# show security screen ids-option untrust-screen
icmp \{
    ping-death;
\}
ip \{
    source-route-option;
    tear-drop;
\}
tcp \{    Note: UDP flood setting also exist
    syn-flood \{
        alarm-threshold 10024;
        attack-threshold 200;
        source-threshold 10024;
        destination-threshold 2048;
        timeout 20;
    \}
    land;
\}
\end{alltt}



\slidenext{}


\end{document}
